\section{Screening vs.~Forensic Stages: An Architectural Test}
\label{sec:forensic}

The contribution claim of §\ref{sec:intro} is structural: under
incomplete observability, enforcement architecture separates an
award-layer screening stage that produces candidate flags from a
bid-layer forensic stage that adjudicates suspected coordination
on the bid distribution. The two stages answer different
questions and require different statistics. This section tests
the architectural claim empirically: the two stages should
discriminate the same cartel-adjacency labels at comparable
levels on different data envelopes, and the combination should
beat either component alone.

\paragraph{Comparable accuracy on different data envelopes}
Table~\ref{tab:imhof_full} reports the comparison on the full
always-loser pool ($N = 16{,}779$ firms, $\valCobidders$
adjudicated cobidder labels). The screening statistic
discriminates at AUC $0.903$ ($95\%$ CI $[0.884, 0.923]$) under
the binary FL14 indicator and $0.884$ ($[0.860, 0.908]$) under
the continuous $\log(1+\text{tenders\_count})$ score. The full
Imhof--Wallimann pipeline---five within-tender bid-distribution
features (coefficient of variation, skewness, kurtosis, spread,
second-low ratio)
\citep{imhof2019detecting,wallimann2023machine}---reaches AUC
$0.888$ ($[0.865, 0.911]$) on the same target. The two pipelines
reach statistically indistinguishable discrimination on
different data envelopes. The construct does not outperform the
bid-distribution pipeline, and we do not claim it does; the
architectural point is that comparable accuracy is achievable
on a substantially thinner data envelope.

% comparison replaces the v13 strawman (CV-only) with the actual five-feature
% pipeline. Source: scripts/31_imhof_full_pipeline.R.
\begin{table}[htbp]
\hfuzz=300pt\relax
\centering
\caption{Full Imhof--Wallimann Pipeline vs Loser-Side Concentration: Discrimination at Different Data Costs}
\label{tab:imhof_full}
\resizebox{\textwidth}{!}{%
\begin{threeparttable}
\small
\begin{tabular}{lccc}
\toprule
Model & Features required & AUC & 95\% CI \\
\midrule
\multicolumn{4}{l}{\emph{Participation-only (award-record level)}} \\
FL flag (binary)            & participation count $>$ 14, win flag                   & $0.903$ & $[0.884, 0.923]$ \\
$\log(1{+}\text{tenders\_count})$ & participation count, win flag                    & $0.884$ & $[0.860, 0.908]$ \\
\midrule
\multicolumn{4}{l}{\emph{Bid-distribution (per-bid microdata)}} \\
Imhof CV only (v13 strawman) & per-bid offer values                                  & $0.585$ & $[0.553, 0.616]$ \\
Imhof full pipeline          & 5 within-tender features (CV, skew, kurt, spread, second-low) & $0.888$ & $[0.865, 0.911]$ \\
\midrule
\multicolumn{4}{l}{\emph{Combined}} \\
FL flag $+$ Imhof full       & participation count + 5 bid features                  & $0.955$ & $[0.943, 0.967]$ \\
$\log(1{+}\text{tenders\_count})$ $+$ Imhof full & participation count + 5 bid features & $0.962$ & $[0.954, 0.969]$ \\
\bottomrule
\end{tabular}
\begin{tablenotes}
\small
\item \textit{Notes:} 5-fold cross-validated AUC against $193$
CADE-cobidder labels within the always-loser pool ($N\!=\!16{,}779$).
Random forests with $500$ trees, otherwise default tuning. Bid-level
features are aggregated to the firm level by averaging the per-tender
within-tender statistic across each firm's tenders.
The two top blocks compare alternative single-construct screens
operating on different data envelopes; the bottom block reports the
combined classifiers. The combined model gains $5$--$7$ AUC points
over either single construct, indicating non-redundant signal in the
participation-only and bid-distribution information sources. The
participation-only screen achieves discrimination comparable to the
full Imhof--Wallimann pipeline at substantially lower data cost,
which is the operational claim of the paper. We do not claim the
participation-only screen \emph{outperforms} the bid-distribution
pipeline; the relevant comparison is at fixed data cost, where the
two pipelines reach similar accuracy and the combination dominates.
\end{tablenotes}
\end{threeparttable}
}
\end{table}


\paragraph{Non-redundant signal} The combined specifications in
the bottom panel of Table~\ref{tab:imhof_full} address the
integration question. FL14 combined with the full Imhof pipeline
reaches AUC $0.955$ ($[0.943, 0.967]$); continuous loss
intensity combined with the full Imhof pipeline reaches AUC
$0.962$ ($[0.954, 0.969]$). The combined classifiers gain
$5$--$7$ AUC points over either single construct. Restricting
attention to the same-sample audit
(Table~\ref{tab:imhof_incremental}), where all detectors
evaluate the identical subset of firms with bid microdata
available, the screening statistic adds $+0.035$ AUC over the
Imhof full pipeline alone (DeLong $p = 0.014$), and the combined
model adds $+0.096$--$+0.098$ AUC over Imhof full alone
($p < 0.001$). The award-layer screening object and the
bid-layer forensic moments carry non-redundant information
about the same cartel-adjacency target.

\paragraph{The architectural reading}
The two findings together identify the sequencing structure
the contribution claims. The screening statistic and the
bid-distribution detectors discriminate the same labels at
comparable levels on different data envelopes (the architecture
treats them as substitutable when only one envelope is
available); the two carry non-redundant signal when both
envelopes are available (the architecture treats them as
sequential when both are). Under incomplete observability the
deployable architecture is therefore: an award-layer screening
stage that runs on contract-record data and produces candidate
environments, followed by a bid-layer forensic stage that runs
the bid-distribution moment battery on the candidates the
screening stage flags, where bid microdata are recoverable for
those candidates. The two stages answer different questions and
require different statistics; the comparison in this section
shows that they share the same cartel-adjacency target without
dominating each other.

The implication for portability is concrete. Jurisdictions that
preserve only the award layer can deploy the screening stage
immediately, on data they already maintain. Jurisdictions that
also preserve bid microdata gain the forensic stage on top, with
the screening stage already constraining the pool of cases that
warrant the forensic step. Reforms that mandate bid-microdata
archival (\S\ref{sec:institutional}) are best read as expanding
the forensic stage rather than substituting for the screening
one---a reading that reverses the standard policy intuition that
bid-archival is a precondition for proactive cartel screening.
